% preface.tex — Preface for Fixed Point Theory course
\chapter*{Preface}
\addcontentsline{toc}{chapter}{Preface}

Fixed point theory stands as one of the fundamental pillars of modern
mathematical analysis. Its influence extends well beyond the boundaries of
functional analysis, permeating such diverse fields as topology, algebra,
optimization, game theory, mathematical economics, and partial differential
equations.

\section*{Objectives}

This course has been designed for Master's and PhD students in mathematics.
It aims to present in a rigorous and thorough manner the major fixed point
theorems, their proofs, their generalizations, and their applications.

The main pedagogical objectives are:
\begin{itemize}
    \item Master the complete proofs of the classical theorems (Banach, Brouwer,
          Schauder, Kakutani, Tarski--Knaster);
    \item Understand the deep connections between these results and the underlying
          mathematical structures (metric, topological, order-theoretic);
    \item Apply these theorems to the resolution of concrete problems:
          differential equations, integral equations, Nash equilibria;
    \item Explore modern extensions: generalized metric spaces, random fixed
          points, multivalued correspondences.
\end{itemize}

\section*{Organization}

The course is structured in eleven chapters organized into three natural parts.

\subsection*{Part I: Metric Theorems (Chapters 1--4)}

The first part is devoted to metric-type results. After an introductory chapter
that sets the historical motivations and fundamental definitions, Chapter~2
presents the celebrated Banach Contraction Principle (1922) and its direct
generalizations. Chapter~3 explores metric extensions due to Edelstein, Kannan,
\'Ciri\'c, and Caristi, which relax the contraction hypothesis in various ways.
Chapter~4 introduces generalized metric spaces ($b$-metrics, $G$-metrics) and
examines how the classical theorems extend to these settings.

\subsection*{Part II: Topological Theorems (Chapters 5--7)}

The second part addresses theorems of a topological nature. Chapter~5 is
dedicated to Brouwer's theorem (1911), a true jewel of algebraic topology,
with several proofs (combinatorial via Sperner's lemma, analytical,
homological). Chapter~6 treats Schauder's theorem (1930), the
infinite-dimensional generalization of Brouwer's theorem, and its variants
(Schauder--Tychonoff, Darbo). Chapter~7 presents Kakutani's theorem (1941) for
multivalued correspondences, an indispensable tool in game theory and
mathematical economics.

\subsection*{Part III: Extensions and Applications (Chapters 8--11)}

The third part gathers advanced topics. Chapter~8 studies fixed points in
ordered spaces (theorems of Tarski, Knaster--Tarski, Bourbaki--Witt). Chapter~9
is a cross-cutting applications chapter: ordinary differential equations,
partial differential equations, Nash equilibria, economic models. Chapter~10
presents the Markov--Kakutani theorem and its consequences for common fixed
points of families of transformations. Finally, Chapter~11 introduces random
fixed point theory, a rapidly developing area.

\section*{Prerequisites}

The reader should possess solid knowledge of:
\begin{itemize}
    \item \textbf{Functional analysis}: Banach spaces, Hilbert spaces, bounded
          linear operators, the Hahn--Banach theorem, weak compactness;
    \item \textbf{General topology}: compact spaces, connectedness, separation
          theorems, complete metric spaces;
    \item \textbf{Measure theory}: measurable spaces, measurable functions,
          Lebesgue integration (for Chapter~11);
    \item \textbf{Linear algebra}: finite-dimensional vector spaces, determinants,
          multilinear forms.
\end{itemize}

\section*{Notation}

Throughout this course, we use the following notation:
\begin{itemize}
    \item $\R$, $\N$, $\Z$, $\Q$, $\C$: the sets of real, natural, integer,
          rational, and complex numbers;
    \item $(X, d)$: a metric space, where $d$ denotes the metric;
    \item $\norm{\cdot}$: a norm on a normed vector space;
    \item $\inner{\cdot}{\cdot}$: an inner product on a Hilbert space;
    \item $B(x, r)$: the open ball with center $x$ and radius $r$;
    \item $\overline{B}(x, r)$: the closed ball with center $x$ and radius $r$;
    \item $\overline{A}$: the closure of a set $A$;
    \item $\mathrm{co}(A)$, $\overline{\mathrm{co}}(A)$: the convex hull and
          closed convex hull of $A$;
    \item $\mathrm{Fix}(T)$: the set of fixed points of the mapping $T$;
    \item $\mathrm{Lip}(T)$: the Lipschitz constant of $T$.
\end{itemize}

\section*{Conventions}

Unless otherwise stated:
\begin{itemize}
    \item Vector spaces are over the field $\R$ of real numbers;
    \item Topological spaces are assumed to be Hausdorff;
    \item Maps take values in $\R$ or in a Banach space;
    \item The symbol $\square$ marks the end of a proof;
    \item The symbol $\lozenge$ marks the end of an example or remark.
\end{itemize}

\section*{Reading guide}

This course may be approached in several ways depending on the reader's level
and interests:

\begin{itemize}
    \item \textbf{Foundational track} (M1): Chapters 1, 2, 5, and 9. This track
          covers the Banach contraction principle and Brouwer's theorem with
          their main applications.
    \item \textbf{Advanced track} (M2): Add Chapters 3, 6, 7, and 8. This
          completes the training with metric extensions, Schauder's theorem,
          Kakutani's theorem, and ordered spaces.
    \item \textbf{Research track} (PhD): The entire course, with particular
          attention to Chapters 4, 10, and 11, which touch on current research
          themes.
\end{itemize}

\section*{Acknowledgments}

This course owes much to the classical references in the field, notably the
works of Granas and Dugundji \cite{granas2003}, Goebel and Kirk
\cite{goebel1990}, Zeidler \cite{zeidler1986}, Agarwal, Meehan, and O'Regan
\cite{agarwal2001}, and Smart \cite{smart1974}. We also thank the many
colleagues and students whose remarks have helped improve this text.

\section*{How to use this course}

Each chapter contains:
\begin{itemize}
    \item Numbered, precise, and motivated \textbf{definitions};
    \item \textbf{Theorems} with complete proofs;
    \item Illustrative \textbf{examples} and \textbf{counter-examples} showing
          the necessity of hypotheses;
    \item \textbf{Remarks} illuminating subtle points;
    \item \textbf{Exercises} of varied difficulty, some with hints.
\end{itemize}

Special boxes (key formulas, warnings, intuition, algorithms) allow quick
identification of essential information.

\bigskip
\hfill \textit{The author}

\hfill \textit{March 2026}
